\documentclass{beamer}

\usepackage[T2A]{fontenc}
\usepackage[utf8]{inputenc}
\usepackage[english,russian]{babel}
\input{settings.tex}

\newcommand{\myqrframe}{
	\begin{frame}
		\centerline{Слайды доступны на Github:}
		\bigskip
		\centerline{\mygit}
		\bigskip
		\myqr
	\end{frame}
}


\subtitle{Современные методы управления роботами}
\author{Сергей Савин}
\date{2026}

\title{Анализ устойчивости методом суммы квадратов}
\centering



\begin{document}
\maketitle

\begin{frame}{Содержание}
	\begin{itemize}
		\item Устойчивость как задача SoS
		\item Асимптотическая устойчивость как задача SoS
		\item Экспоненциальная устойчивость как задача SoS
		\item Квадратичная функция Ляпунова
		\item Неполиномиальные системы
		\item Проверка устойчивости маятника
		\item Локальная устойчивость
		\item S-процедура
		\item Локальная асимптотическая устойчивость с помощью S-процедуры
		\item Литература
	\end{itemize}
\end{frame}

\begin{frame}{Устойчивость как задача SoS}
	\begin{flushleft}
		Система $\dot{x}=f(x)$ является устойчивой, если существует $V(x)$ такая, что:
		\[
		V(x)>0,\qquad V(0)=0,\qquad \dot V(x)\leq0\footnotemark
		\]
		
		\textbf{Замечание:}
		\begin{itemize}
			\item Строгое неравенство заменяется ограничением:
			\[
			V(x)\geq\alpha x^\top x,\qquad \alpha>0.
			\]
			
			\item Поскольку $\dot V(x)
			=
			\frac{\partial V}{\partial x}f(x),$	можно записать $-\frac{\partial V}{\partial x}f(x)\geq0$.
		\end{itemize}
		
		Таким образом, критерий устойчивости можно сформулировать как SoS-задачу:
		\[
		\begin{array}{ll}
			\text{find} & V\\[0.2em]
			\text{subject to:} &
			V(x)-\alpha x^\top x\in\mathrm{SoS},\\
			&
			-\dfrac{\partial V}{\partial x}f(x)\in\mathrm{SoS}.
		\end{array}
		\]
	\end{flushleft}
	
	\footnotetext{\vspace{1.0em}
		Условие $V(x)\geq\alpha x^\top x$, $\alpha>0$ означает, что $V(x)\to\infty$ при $\|x\|\to\infty$.}
\end{frame}


\begin{frame}{Асимптотическая устойчивость как задача SoS}
	\begin{flushleft}
		Система $\dot{x}=f(x)$	является асимптотически устойчивой, если существует $V(x)$ такая, что:
		\[
		V(x)>0,\qquad V(0)=0,\qquad
		\textcolor{myblue}{\dot V(x)<0}.
		\]
		
		Строгое неравенство заменяется условием
		\[
		\dot V(x)\leq-\epsilon x^\top x,
		\qquad \epsilon>0.
		\]
		
		Получаемая SoS-задача имеет вид:
		\[
		\begin{array}{ll}
			\text{find} & V\\[0.2em]
			\text{subject to:} &
			V(x)-\alpha x^\top x\in\mathrm{SoS},\\
			&
			-\dfrac{\partial V}{\partial x}f(x)
			\textcolor{myblue}{-\epsilon x^\top x}\in\mathrm{SoS}.
		\end{array}
		\]
	\end{flushleft}
	
\end{frame}


\begin{frame}{Экспоненциальная устойчивость как задача SoS}
	\begin{flushleft}
		Система $\dot{x}=f(x)$	является экспоненциально устойчивой, если существуют $V(x)$ и
		$\alpha,\beta,\epsilon>0$ такие, что:
		\[
		\alpha x^\top x
		\leq V(x)
		\textcolor{myblue}{\leq \beta x^\top x},\qquad 
		\dot V(x)\leq-\epsilon x^\top x.
		\]
		
		\bigskip
		
		Соответствующая SoS-задача имеет вид:
		\[
		\begin{array}{ll}
			\text{find} & V\\[0.2em]
			\text{subject to:} &
			V(x)-\alpha x^\top x\in\mathrm{SoS},\\
			&
			\textcolor{myblue}{-V(x)+\beta x^\top x\in\mathrm{SoS}},\\
			&
			-\dfrac{\partial V}{\partial x}f(x)
			-\epsilon x^\top x\in\mathrm{SoS}.
		\end{array}
		\]
	\end{flushleft}
	
\end{frame}


\begin{frame}{Квадратичная функция Ляпунова}
	\begin{flushleft}
		Можно упростить задачу, ограничившись квадратичными функциями Ляпунова:
		\[
		V(x)=\frac{1}{2}x^\top Px,
		\qquad P=P^\top.
		\]
		
		Например, критерий устойчивости принимает вид:
		\[
		\begin{array}{ll}
			\text{find} & P\\[0.2em]
			\text{subject to:} &
			P\succeq\alpha I,\\
			&
			-x^\top Pf(x)\in\mathrm{SoS}.
		\end{array}
		\]
		
		
		Заменяя последнее условие на:
		\[
		-x^\top Pf(x)-\epsilon x^\top x\in\mathrm{SoS}
		\]
		мы проверяем асимптотическую и экспоненциальную устойчивость. Действительно, для любой $P\succ0$:
		\[
		\lambda_{\max}(P)x^\top x
		\geq V(x)
		\geq
		\lambda_{\min}(P)x^\top x.
		\]
	\end{flushleft}
\end{frame}

\begin{frame}{Неполиномиальные системы}
	\begin{flushleft}
		Многие динамические системы являются неполиномиальными.
		Например, механические системы часто содержат тригонометрические функции.
		
		\bigskip
		
		Рассмотрим маятник:
		\[
		ml^2\ddot{\theta}
		+b\dot{\theta}
		+mgl\sin(\theta)=0.
		\]
		
		\bigskip
		
		Работа с такими системами требует полиномиальной переформулировки:
		
		\begin{itemize}
			\item Аппроксимировать нелинейности, например:
			\[
			\sin(x)\approx
			x-\frac{x^3}{3!}+\frac{x^5}{5!}.
			\]
			
			\item Найти эквивалентную полиномиальную формулировку,
			возможно, используя дополнительные переменные и ограничения:
			\[
			s=\sin(x),\qquad
			c=\cos(x),\qquad
			s^2+c^2=1.
			\]
		\end{itemize}
	\end{flushleft}
	
\end{frame}


\begin{frame}{Проверка устойчивости маятника, часть 1}
	\begin{flushleft}
		Рассмотрим простой маятник, у которого все параметры равны $1$:
		\[
		\ddot{\theta}+\dot{\theta}+\sin(\theta)=0.
		\]
		
		\bigskip
		
		Введём новые переменные:
		\[
		x_1=\dot{\theta},
		\qquad
		x_2=\sin(\theta),
		\qquad
		x_3=\cos(\theta)-1.
		\]
		
		\bigskip
		
		Эта замена переменных позволяет сделать четыре вещи:
		
		\begin{enumerate}
			\item Преобразовать дифференциальное уравнение второго порядка в систему первого порядка.
			
			\item Исключить тригонометрические функции из динамики.
			
			\item Перевести интересующее нас положение равновесия в точку $x=0$.
			
			\item Исключить $\theta$ из числа переменных состояния, поэтому нам не нужно устанавливать связь между $\theta$ и $x_2$, $x_3$.
		\end{enumerate}
	\end{flushleft}
	
\end{frame}


\begin{frame}{Проверка устойчивости маятника, часть 2}
	\begin{flushleft}
		В новых переменных динамика $\ddot{\theta}+\dot{\theta}+\sin(\theta)=0$ принимает вид:
		\[
		\begin{aligned}
			\dot{x}_1 &= -x_1-x_2, \\
			\dot{x}_2 &= \cos(\theta)\dot{\theta} = x_1(x_3+1), \qquad  \textcolor{mygray}{\cos(\theta)=x_3+1} \\
			\dot{x}_3 &= -\sin(\theta)\dot{\theta} = -x_1x_2,  \qquad \quad  \textcolor{mygray}{\sin(\theta)=x_2}
		\end{aligned}
		\]
		
		
		\bigskip
		
		Связь между синусом и косинусом задаётся соотношением
		\[
		\sin^2(\theta)+\cos^2(\theta)=1.
		\]
		
		Используя новые переменные, это соотношение принимает вид
		\[
		h(x)
		=
		x_2^2+(x_3+1)^2-1
		=0.
		\]
		
		\bigskip
		
		Таким образом, мы получаем полиномиальную систему с полиномиальным ограничением:
		\[
		\boxed{
			\dot{x}=f(x),
			\qquad
			h(x)=0.
		}
		\]
	\end{flushleft}
\end{frame}


\begin{frame}{Проверка устойчивости маятника, часть 3}
	\begin{flushleft}
		Проверку устойчивости можно сформулировать с помощью SoS-программирования. Важно, что условия устойчивости должны выполняться только на многообразии
		\[
		h(x)=0.
		\]
		
		
		Для этого можно использовать полиномиальные множители:
		\[
		\begin{array}{ll}
			\text{find} & V(x),\ \lambda(x),\ \mu(x)\\[0.2em]
			\text{subject to:}
			&
			V(x)-\alpha x^\top x+\mu(x)h(x)
			\in\mathrm{SoS},
			\\[0.4em]
			&
			-\dfrac{\partial V}{\partial x}f(x)
			-\epsilon x^\top x+\lambda(x)h(x)
			\in\mathrm{SoS}.
		\end{array}
		\]
		
		\bigskip
		
		Заметим, что $\lambda(x)$ и $\mu(x)$ являются произвольными полиномиальными множителями;
		они не обязаны быть неотрицательными. На многообразии $h(x)=0$ эти условия сводятся к
		\[
		V(x)\geq\alpha x^\top x,
		\qquad
		\dot V(x)\leq-\epsilon x^\top x.
		\]
		
		
	\end{flushleft}
\end{frame}


\begin{frame}{Локальная устойчивость}
	\begin{flushleft}
		Многие системы, включая маятник, устойчивы только локально. Для проверки локальной устойчивости потребуем, чтобы условия Ляпунова выполнялись только в некоторой области $\mathcal{D}$.
		
		\bigskip
		
		Опишем $\mathcal{D}$ как множество уровня положительно определённой функции $W(x)$:
		\[
		\mathcal{D}
		=
		\left\{
		x: W(x)\leq c
		\right\},
		\qquad c>0.
		\]
		
		\bigskip
		
		Например, если
		\[
		W(x)=x^\top Px,
		\qquad P\succ0,
		\]
		то
		\[
		\mathcal{D}
		=
		\left\{
		x:x^\top Px\leq c
		\right\}
		\]
		является эллипсоидом.
	\end{flushleft}
\end{frame}


\begin{frame}{S-процедура}
	\begin{flushleft}
		Предположим, что мы хотим доказать, что полином $f(x)$ неотрицателен на множестве
		\[
		g(x)\leq0.
		\]
		
		Мы можем искать полиномиальный множитель
		$\lambda(x)\in\mathrm{SoS}$ такой, что:
		
		\begin{align}
			\label{eq:s-procedure}
			f(x)+\lambda(x)g(x)\in\mathrm{SoS}.
		\end{align}
		
		
		\begin{itemize}
			\item Уравнение \eqref{eq:s-procedure} означает $
			f(x)+\lambda(x)g(x)\geq0$. Для таких $x$, что $g(x)\leq0$, имеем $f(x)
			\geq
			-\lambda(x)g(x)
			\geq0.$
			
			\item Однако для таких $x$, что $g(x)>0$, имеем $-\lambda(x)g(x)\leq0$, 
			поэтому данное условие не означает, что $f(x)\geq0$.
		\end{itemize}
		
		
		\bigskip
		
		Это позволяет SoS-задаче обеспечивать выполнение условия $f(x)\geq0$
		только на требуемом множестве $g(x)\leq0$.
	\end{flushleft}
\end{frame}




\begin{frame}{Локальная асимптотическая устойчивость с помощью S-процедуры}
	\begin{flushleft}
		С помощью S-процедуры можно ограничить требование положительности
		$V(x)$ и требование отрицательности $\dot V(x)$ некоторой областью $\mathcal{D}$. Это позволяет доказать \textbf{локальную асимптотическую устойчивость}.
		
		\bigskip
		
		Естественным выбором $\mathcal{D}$ является подуровень $V(x)$:
		\[
		\mathcal{D}
		=
		\left\{
		x:V(x)\leq c
		\right\}.
		\]
		
		\bigskip
		
		SoS-формулировка принимает вид
		\[
		\begin{array}{ll}
			\text{find} &
			V,\ \lambda_1,\ \lambda_2
			\\[0.2em]
			\text{subject to:}
			&
			V(x)-\alpha x^\top x
			+\lambda_1(x)\bigl(V(x)-c\bigr)
			\in\mathrm{SoS},
			\\[0.4em]
			&
			-\dfrac{\partial V}{\partial x}f(x)
			-\epsilon x^\top x
			+\lambda_2(x)\bigl(V(x)-c\bigr)
			\in\mathrm{SoS},
			\\[0.4em]
			&
			\lambda_1(x)\in\mathrm{SoS},
			\qquad
			\lambda_2(x)\in\mathrm{SoS}.
		\end{array}
		\]
		{\small
			$\lambda_1$ и $\lambda_2$ должны быть SoS
			(неотрицательными) множителями. Параметры $\alpha>0$
			и $\epsilon>0$ обеспечивают выполнение строгих неравенств Ляпунова.}
	\end{flushleft}
\end{frame}




\begin{frame}{Литература}
	
	\begin{itemize}
		
		\item Stephen Lall -- Stanford University:
		\bref{https://lall.stanford.edu/ee464/lectures/sum_of_squares.pdf}
		{Sum of Squares.}
		
		\item Russ Tedrake -- MIT:
		\bref{https://underactuated.csail.mit.edu/lyapunov.html}
		{Underactuated Robotics: Lyapunov Analysis, ch. 9.2.2, 9.2.6}
		
		\item MIT OpenCourseWare -- Algebraic Techniques and Semidefinite Optimization:
		\bref{https://ocw.mit.edu/courses/6-972-algebraic-techniques-and-semidefinite-optimization-spring-2006/resources/lecture_10/}
		{Lecture 10: Sums of Squares and Semidefinite Programming.}
		
		\item Antonis Papachristodoulou and Stephen Prajna:
		\bref{https://users.ox.ac.uk/~engs0587/PapPACC05.pdf}
		{A Tutorial on Sum of Squares Techniques for Systems Analysis.}
		
		
	\end{itemize}
	
\end{frame}


\myqrframe


\end{document}
